\scanpage{17}
\begin{center}
INSTITUT\\
DES HAUTES ÉTUDES SCIENTIFIQUES\\
\rule{2.3em}{0.4pt}
\end{center}

\uline{Continuation de} SGA 4 $\underline{\overline{\mathrm{IV}}}$ 4.10.6 (\makebox[0pt][l]{\tikz[overlay,baseline=0pt,x=1em,y=1em]{\draw[gray,thin] (-.45,-1.1) .. controls (-.55,-.25) and (.1,-.35) .. (1.3,-.35) -- (8.1,-.35) .. controls (8.9,-.35) and (9.1,-.1) .. (9,.55);}}\inserted{généralisation des} 3° et 4°\unsure{)} appliqué à $S(x)\subset S_{/x}$)

5°) Soit $S_0\hookrightarrow S$ une \deleted{\unsure{foncteur|foncteurs} ($S_0$ cat. \unsure{$\in\mathcal{U}$|$\subset\mathcal{U}$|$\mathcal{U}$}) et supposons que} sous-catégorie pleine de $S$ satisfaisant les conditions suivantes :

(i) Si $u:X\mathbin{-}Y$ avec $u\in\underline{M}$, $Y\in\mathrm{ob}\,S_0$, donc $X\in\mathrm{ob}\,S_0$

[ (ii) \deleted{$\forall X\in\mathrm{ob}\,S$, existe famille \unsure{couvrante|couvrant} $X_i\to X$, \unsure{telle} \unsure{que $X_i$|$Y_i\mathbin{-}X_i$} \unsure{tq|soit}} le crible de $S$ engendré par les objets de $S_0$ est couvrant, i.e. $\forall X\in\mathrm{ob}\,S$, existe famille couvrante $X_i\to X$, avec $X_i\to Z_i\in\mathrm{ob}\,S_0$. ]

\unsure{.} Munissons $S_0$ de la top. induite. \unsure{Prouver|Prouvons} qu'une famille de morphismes $X_i-X$ dans $S_0$ est couvrante dans $S_0$ ssi elle l'est dans $S_{\deleted{\unsure{0}}}$ et que $S_0\to S$ est \uline{cocontinu.}

Soit $\underline{M}_0$ l'ens des $u\in\underline{M}$, \deleted{\unsure{tq}} $u:x-y$, avec $Y\in\mathrm{ob}\,S_0$, donc $X\in\mathrm{ob}\,S_0$. \unsure{Prouver|Prouvons} que $\underline{M}_0\in\mathrm{Fl}(S_0)$ satisfait les conditions a) b) c) d) de \uline{1}°).

Procédant comme dans 3°, définir des foncteurs
\par\noindent
\begin{tabular}{@{}l@{\quad}l@{}}
$\mathrm{Res}:\widetilde{S}\to\widetilde{S}_0,$ & \makebox[0pt][l]{\tikz[overlay,baseline=0pt,x=1em,y=1em]{\draw[gray,thin] (9.4,1.35) .. controls (10,2.1) and (-.9,1.7) .. (-.8,.35) .. controls (-1,-2.3) and (8.8,-2.3) .. (10,-.65) .. controls (10.5,.15) and (10.1,.8) .. (9.8,1.05);}}$\mathrm{Prol}:\widetilde{S}_0\mathbin{-}\widetilde{S}_{\deleted{\unsure{0|2}}}\;?$\\
\deleted{et \unsure{prouver|prouvons} que}\inserted{tels que} $\mathrm{Res}$ comm. & aux limites
\end{tabular}\par\noindent
inductives et projectives, et $\mathrm{Prol}$ est exact \unsure{à|:} g. et pl fidèle, son image essentielle
% archive sheet number: 16
